\section{Arrow diagrams}

A complete lattice is a category with all (co)limits: objects are elements, there is one
arrow $a \to b$ exactly when $a \le b$, $\Sup$ is the colimit and $\Inf$ the limit.  Under
that dictionary each lemma of the development becomes a diagram, and each bridge of
\lean{RequestProject/Bridges.lean} becomes a functor.

\subsection{The proofs as commuting triangles}

Both collapse lemmas are a \lean{trans} of two equalities, i.e.\ a commuting triangle
whose slanted edge is the theorem.

\begin{figure}[ht]
\centering
\begin{minipage}{.48\textwidth}\centering
\begin{cdscale}{\textwidth}
\begin{tikzcd}[column sep=2.6cm,row sep=1.7cm]
\displaystyle\Sup_i \Sup_j F_{ij}
  \arrow[r,"\lab{iSup\_pprod}"]
  \arrow[dr,swap,pos=.42,"\lab{iSup₂\_eq\_iSup\_diagonal}"]
& \displaystyle\Sup_{p:\iota\times\iota} F_{p_1p_2}
  \arrow[d,"\lab{iSup\_eq\_iSup\_of\_cofinal}"] \\
& \displaystyle\Sup_k F_{kk}
\end{tikzcd}
\end{cdscale}
\\[4pt]{\footnotesize complete lattice: no side conditions}
\end{minipage}\hfill
\begin{minipage}{.48\textwidth}\centering
\begin{cdscale}{\textwidth}
\begin{tikzcd}[column sep=2.6cm,row sep=1.7cm]
\displaystyle\Sup_i \Sup_j F_{ij}
  \arrow[r,"\lab{ciSup\_pprod}"]
  \arrow[dr,swap,pos=.42,"\lab{ciSup₂\_eq\_ciSup\_diagonal}"]
& \displaystyle\Sup_{p:\iota\times\iota} F_{p_1p_2}
  \arrow[d,"\lab{ciSup\_eq\_ciSup\_of\_cofinal}"] \\
& \displaystyle\Sup_k F_{kk}
\end{tikzcd}
\end{cdscale}
\\[4pt]{\footnotesize conditionally complete: the horizontal arrow, and only it, consumes
\lean{bddAbove\_range\_uncurry}}
\end{minipage}
\caption{The same triangle twice.  Passing from the complete to the conditionally
complete world changes one edge label, not the shape.}
\end{figure}

\subsection{The binary-operation form as a zig-zag}

\lean{Golf.iSup\_op\_iSup\_diagonal} takes the two distributivity laws of the ambient
theory as parameters and produces the diagonal form; the proof is a three-step chain.

\begin{figure}[ht]
\centering
\begin{tikzcd}[column sep=2.4cm,row sep=1.4cm]
\mathrm{op}\bigl(\Sup_i f_i\bigr)\bigl(\Sup_j g_j\bigr)
  \arrow[r,"h_\ell"]
& \displaystyle\Sup_i \mathrm{op}\,(f_i)\bigl(\Sup_j g_j\bigr)
  \arrow[d,"\Sup_i h_r"] \\
\displaystyle\Sup_k \mathrm{op}\,(f_k)(g_k)
& \displaystyle\Sup_i\Sup_j \mathrm{op}\,(f_i)(g_j)
  \arrow[l,swap,"\lab{iSup₂\_eq\_iSup\_diagonal}"]
\end{tikzcd}
\caption{Instantiating $\mathrm{op}$ with \lean{(·+·)} and the pair
\lean{ENNReal.iSup\_add}/\lean{ENNReal.add\_iSup} gives \lean{ennreal\_iSup\_add\_iSup};
with \lean{(·*·)} and \lean{ENNReal.iSup\_mul}/\lean{ENNReal.mul\_iSup} the multiplicative
analogue --- same proof, new inputs.}
\end{figure}

\subsection{Currying and Fubini for colimits}

\lean{Golf.iSup\_pprod} is the lattice shadow of the Fubini theorem for colimits: the
colimit over a product category is the iterated colimit.

\begin{figure}[ht]
\centering
\begin{tikzcd}[column sep=1.8cm,row sep=1.3cm]
\iota\times\iota
  \arrow[r,"F"]
  \arrow[d,swap,"\mathrm{pr}_1"]
& \alpha \\
\iota
  \arrow[ur,swap,"{\Sup_j F(-,j)}"]
&
\end{tikzcd}
\qquad\qquad
\begin{tikzcd}[column sep=1.6cm,row sep=1.3cm]
\iota
  \arrow[r,"\Delta"]
  \arrow[dr,swap,"F\circ\Delta"]
& \iota\times\iota
  \arrow[d,"F"] \\
& \alpha
\end{tikzcd}
\caption{Left: the left Kan extension along $\mathrm{pr}_1$ is the row-supremum, and its
colimit is $\Sup_i\Sup_j F_{ij}$.  Right: the diagonal functor
$\Delta = $ \lean{CategoryTheory.Functor.diag}, whose composite with $F$ is the diagonal
family.  The theorem says the two colimits agree.}
\end{figure}

\subsection{The categorical bridge: the collapse \emph{is} finality of $\Delta$}

For a bimonotone family over a directed (= filtered) index, no order-theoretic argument
is needed at all: $\Delta$ is a final functor, and finality is precisely invariance of
colimits: this is
\lean{iSup\_prod\_eq\_iSup\_diagonal\_functor} in \lean{Bridges.lean}.

\begin{figure}[ht]
\centering
\begin{tikzcd}[column sep=2.1cm,row sep=1.35cm]
\iota
  \arrow[r,"\Delta"]
  \arrow[dr,swap,"F\Delta"]
& \iota\times\iota
  \arrow[d,"F"]
  \arrow[r,"\mathrm{colim}"]
& \operatorname{colim} F
  \arrow[d,equals,"\lab{Functor.Final.colimitIso}"{xshift=2mm}] \\
& \alpha
  \arrow[r,swap,"\mathrm{colim}"]
& \operatorname{colim} (F\Delta)
\end{tikzcd}
\qquad
\raisebox{1.1cm}{$\begin{aligned}
\operatorname{colim} F &= \Sup_{p}F_p &&\lab{colimit\_eq\_iSup}\\
\operatorname{colim} F\Delta &= \Sup_k F_{kk} &&\lab{colimit\_eq\_iSup}\\
\Delta \text{ final} &\ \Leftarrow\ \iota \text{ filtered}
  &&\lab{final\_diag\_of\_isFiltered}
\end{aligned}$}
\caption{The categorical route.  Directedness of $\iota$ enters as filteredness of the
index category; cofinality of the diagonal \emph{set} becomes finality of the diagonal
\emph{functor}.}
\end{figure}

\subsection{Transport along the bridges}

Each bridge of \lean{Bridges.lean} is a map of contexts.  What differs between them is
how much of the statement crosses: an order isomorphism carries hypothesis
\emph{and} conclusion, a sup-homomorphism only the conclusion, and duality carries
everything for free because it is the identity on propositions.

\begin{figure}[ht]
\centering
\begin{minipage}{.32\textwidth}\centering
\begin{tikzcd}[column sep=1.1cm,row sep=.95cm]
\alpha \arrow[r,"e","\sim"'] \arrow[d,swap,"\Sup"] & \beta \arrow[d,"\Sup"] \\
\alpha \arrow[r,swap,"e"] & \beta
\end{tikzcd}\\[3pt]
{\footnotesize\lean{OrderIso}: \\ \lean{e.map\_iSup}, both ways}
\end{minipage}\hfill
\begin{minipage}{.32\textwidth}\centering
\begin{tikzcd}[column sep=1.1cm,row sep=.95cm]
\alpha \arrow[r,"e"] \arrow[d,swap,"\Sup"] & \beta \arrow[d,"\Sup"] \\
\alpha \arrow[r,swap,"e"] & \beta
\end{tikzcd}\\[3pt]
{\footnotesize\lean{sSupHom}: \emph{not} injective; \\ only the conclusion crosses}
\end{minipage}\hfill
\begin{minipage}{.32\textwidth}\centering
\begin{tikzcd}[column sep=1.25cm,row sep=.95cm]
\alpha
  \arrow[r,bend left=25,"l"]
  \arrow[r,phantom,"\scriptstyle\bot"]
  \arrow[d,swap,"\Sup"]
& \beta
  \arrow[l,bend left=25,"u"]
  \arrow[d,"\Sup"] \\
\alpha \arrow[r,swap,"l"] & \beta
\end{tikzcd}\\[3pt]
{\footnotesize\lean{GaloisConnection}: the left \\ adjoint preserves $\Sup$}
\end{minipage}
\caption{Three functorial bridges of increasing weakness.  In all three the square
commutes, which is all the proof of the transported collapse uses
(\lean{simp only [← map\_iSup, iSup₂\_eq\_iSup\_diagonal]}).}
\end{figure}

\begin{figure}[ht]
\centering
\begin{tikzcd}[column sep=2.4cm,row sep=1.2cm]
\alpha
  \arrow[r,"{(-)^{\mathrm{od}}}","\sim"']
  \arrow[d,swap,"{\Sup_i\Sup_j F_{ij} = \Sup_k F_{kk}}"]
& \alpha^{\mathrm{od}}
  \arrow[d,"{\Inf_i\Inf_j F_{ij} = \Inf_k F_{kk}}"] \\
\mathrm{Prop}
  \arrow[r,equals,"\lab{Iff.rfl}"]
& \mathrm{Prop}
\end{tikzcd}
\caption{The duality bridge (\lean{Golf.Bridge.iSup₂\_diagonal\_iff\_dual}).  The
$\Inf$-statement in $\alpha^{\mathrm{od}}$ and the $\Sup$-statement in $\alpha$ are not
merely equivalent: they are the \emph{same} proposition, which is why every dual lemma
in \lean{Diagonal.lean} is proved by its primal with the single annotation
\lean{(α := αᵒᵈ)}.}
\end{figure}

\subsection{The bridge hub}

\begin{figure}[ht]
\centering
\begin{tikzpicture}[x=1cm,y=1cm,font=\footnotesize]
  \node[rounded corners=4pt,draw=diagcol,fill=diagcol!8,line width=1pt,
        align=center,inner sep=6pt] (T) at (0,0)
    {\textbf{diagonal collapse}\\[2pt]$\Sup_i\Sup_j F_{ij} = \Sup_k F_{kk}$};
  \tikzset{ctx/.style={rounded corners=3pt,draw=black!45,fill=white,align=center,
    inner sep=3.5pt,font=\scriptsize}}
  \node[ctx] (dual)  at (-5.2, 2.6) {$\alpha^{\mathrm{od}}$\\ duality};
  \node[ctx] (sets)  at (-5.6, 0.0) {$\mathrm{Set}\,\alpha$\\ \lean{IsCofinalFor}};
  \node[ctx] (lin)   at (-5.2,-2.6) {cond.\ complete\\ \emph{linear} order};
  \node[ctx] (iso)   at ( 0.0, 3.0) {$e : \alpha \simeq_o \beta$\\ order iso};
  \node[ctx] (hom)   at ( 5.3, 2.6) {$e : \mathrm{sSupHom}$\\ non-faithful};
  \node[ctx] (gal)   at ( 5.7, 0.0) {$l \dashv u$\\ Galois};
  \node[ctx] (cat)   at ( 5.2,-2.6) {$\mathrm{colim}$, $\Delta$ final\\ category theory};
  \node[ctx] (app)   at ( 0.0,-3.0)
    {$\enn$, $\mathbb{N}_\infty$, $\mathrm{Cardinal}$\\ applications};
  \foreach \n/\lab in {dual/{\lean{Iff.rfl}},sets/{\lean{sSup∘range}},
      lin/{no \lean{BddAbove}},iso/{hyp.\ \& concl.},hom/{concl.\ only},
      gal/{\lean{gc.l\_iSup}},cat/{finality},app/{instances}}{
    \draw[-{Stealth[length=2mm]},draw=bridgecol!80,line width=.7pt]
      (T) -- node[sloped,above,font=\tiny,text=bridgecol]{\lab} (\n);}
\end{tikzpicture}
\caption{One theorem, eight contexts.  The label on each arrow records how much of the
statement survives the transport --- the whole point of the bridge methodology is that
the cheapest context is not always the one the statement was posed in.}
\end{figure}

\clearpage
